\section{Conceptos}

\subsection{Corriente Eléctrica - I}
Es el flujo de electrones que circula a través de un conductor (como un cable).

Unidad: Amperio/Ampere (A)

El ampere se definió como la cantidad de carga eléctrica (coulombs) que pasa por un punto de conductor en un segundo.

1 Ampere = 1 Coulomb/Segundo

\subsection{Voltaje - V}
Es la diferencia de potencial eléctrico entre dos puntos, como la "fuerza" que empuja a los electrones para que circulen.

Unidad: Voltio/Volt (V)

El voltio mide la cantidad de energía (en joules) que tiene cada carga eléctrica (coulomb).

1 Volt = 1 Joule/Coulomb

\subsection{Resistencia - R}
Es la oposición al paso de corriente eléctrica en un material o componente.

Unidad: Ohmio/Ohm ($\Omega$)

Se definió gracias a la Ley de Ohm: $R = \frac{V}{I}$.

\subsection{Ley de Ohm}
Tenemos la Ley de Ohm:
\begin{equation}
    V = RI
\end{equation}
se puede memorizar como "Viva la Reina de Inglaterra".

Y de ahí se derivan:
\begin{equation}
    I = \frac{V}{R}
\end{equation}

\begin{equation}
    R = \frac{V}{I}
\end{equation}

\subsection{Potencia Eléctrica}
También se puede escribir usando la Ley de Ohm:
\begin{equation}
    P = I^{2}R
\end{equation}
o
\begin{equation}
    P = \frac{V^{2}}{R}
\end{equation}

\subsection{Leyes de Resistencias}
En serie:
\begin{equation}
    R = R_1 + R_2 + R_3 + ... + R_n
\end{equation}

En paralelo:
\begin{equation}
    \frac{1}{R} = \frac{1}{R_1} + \frac{1}{R_2} + \frac{1}{R_3} + ... + \frac{1}{R_n}
\end{equation}

Con las fórmulas anteriores, ya se pueden resolver la mayoría de problemas básicos de corriente, voltaje y resistencia.

\subsection{Tip Unidades}
La fórmula de la Ley de Ohm sigue las siguientes unidades:
\begin{equation}
    [V] = [\Omega][A]
\end{equation}

Pero para representar a un Volt, podemos hacer lo siguiente:
\begin{equation}
    [V] = 1000[\Omega][mA] = [K\Omega][mA] = [V]
\end{equation}
donde $1 A = 1000 mA$ y $1 K\Omega = 1000\Omega$.

Así, un Volt se puede representar con las unidades:
\begin{equation}
    [V] = [K\Omega][mA]
\end{equation}

\subsection{Corriente Directa - CC/DC}
La corriente directa es un tipo de flujo eléctrico en el que los electrones se mueven de manera constante en \textbf{una sola dirección}, desde el polo negativo hacia el positivo en un circuito eléctrico cerrado. A diferencia de la corriente alterna, donde el flujo cambia de dirección periódicamente, la corriente directa mantiene una polaridad constante.

Este tipo de corriente es producido por fuentes como baterías, celdas solares y generadores de corriente directa. Se usa principalmente en electrónica y almacenamiento de energía

Extraido de \href{https://www.clark.com.ar/todo-sobre-la-corriente-directa-principios-usos-y-ventajas/}{acá}.

\subsection{Corriente Alterna - AC}
La corriente alterna es un tipo de corriente eléctrica que cambia la polaridad y magnitud en intervalos de tiempo de forma regular. Es decir, \textbf{cambia su dirección} de forma repetida.

Se usa principalmente en electrodomésticos, contactos de casas e iluminación

Extraido de \href{https://universidadeuropea.com/blog/que-es-corriente-alterna/}{acá} y de \href{https://www.clark.com.ar/todo-sobre-la-corriente-directa-principios-usos-y-ventajas/}{acá}.